% ============================================================================
%  3  Stiefel--Whitney 类与欧氏嵌入障碍
% ============================================================================
\section{Stiefel--Whitney 类与欧氏嵌入障碍}
\label{sec:sw}

\subsection{Stiefel--Whitney 类与 Whitney 乘积公式}

对实向量丛 $E\to X$，其总 Stiefel--Whitney 类定义为
\begin{equation*}
  w(E)=1+w_1(E)+w_2(E)+\cdots,
\end{equation*}
其中 $w_i(E)\in H^{i}(X;\Ftwo)$，并约定 $w_0(E)=1$，且当
$i>\operatorname{rank}E$ 时 $w_i(E)=0$。Whitney 乘积公式为
$w(E\oplus F)=w(E)w(F)$；若 $E\oplus F\cong\eps^{r}$，则 $w(F)=w(E)^{-1}$。

本文反复使用下列四条规则。它们都由分裂原理得到：在适当的拉回下把丛拆成线丛的
直和，$w$ 就归约为线丛的一次类。

\begin{enumerate}[label=(\roman*),leftmargin=2.4em,itemsep=0.25em]
  \item \emph{直和}：$w(E\oplus F)=w(E)\,w(F)$，特别地 $w(\eps^{r})=1$。
  \item \emph{张量积}：若 $E$ 的形式根为 $x_1,\dots,x_k$、$F$ 的为
        $y_1,\dots,y_l$，则 $E\otimes F$ 的形式根为 $x_i+y_j$，故
        \begin{equation}
          w(E\otimes F)=\prod_{i=1}^{k}\prod_{j=1}^{l}\bigl(1+x_i+y_j\bigr).
          \label{eq:tensor}
        \end{equation}
        特别地，对线丛 $L,L'$ 有 $w_1(L\otimes L')=w_1(L)+w_1(L')$。
  \item \emph{对偶}：$E^{*}=\Hom(E,\underline{\R})$ 的形式根为 $-x_i$，而在
        $\Ftwo$ 中 $-x_i=x_i$，故 $w_i(E^{*})=w_i(E)$。
  \item \emph{$\Hom$ 丛}：$\Hom(E,F)\cong E^{*}\otimes F$，结合 (ii)(iii) 得
        其形式根同样是 $x_i+y_j$，故
        $w(\Hom(E,F))=w(E\otimes F)$。
\end{enumerate}

我们把 $w(E)^{-1}$ 的次数 $j$ 部分记作 $\wbar_j(E)$，并称 $\wbar(E)$ 为 $E$ 的
\emph{共轭 Stiefel--Whitney 类}。若 $w(E)=1+w_1+w_2+\cdots$，则在 $\Ftwo$ 上
递推得到
\begin{equation*}
  \wbar_0=1,
  \qquad
  \wbar_j=\sum_{i=1}^{j} w_i\,\wbar_{j-i}\quad (j\ge 1).
\end{equation*}
这里的共轭类是向量丛总类的逆，它与正交补丛自身的 Stiefel--Whitney 类不是同一个
对象：后者在第 \ref{sec:classical} 章记作 $v_j$。

\subsection{法丛与切丛的特征类关系}

设 $M^{d}$ 是闭光滑流形，并有光滑嵌入 $f\colon M\emb\R^{d+c}$。管状邻域定理给出
$M$ 上秩为 $c$ 的法丛 $\nu_f$，且
\begin{equation*}
  TM\oplus\nu_f\cong\eps^{d+c},
  \qquad
  w(\nu_f)=w(TM)^{-1}=\wbar(TM).
\end{equation*}
于是，秩为 $c$ 的向量丛在次数 $j>c$ 上自动没有 Stiefel--Whitney 类，即
\begin{equation*}
  \wbar_j(TM)=w_j(\nu_f)=0\quad (j>c).
\end{equation*}
换言之，若 $\wbar_j(TM)\neq0$，则 $c\ge j$。这只是一般的\emph{法丛秩障碍}；
它没有排除 $c=j$ 的临界情形。

\subsection{共轭 Stiefel--Whitney 类给出的嵌入障碍}

\begin{theorem}[法丛秩障碍]\label{thm:embedding-obstruction}
设闭光滑 $d$ 维流形 $M$ 嵌入 $\R^{d+c}$，且 $\wbar_j(TM)\neq0$，则 $c\ge j$。
特别地，若
\begin{equation}
  \delta(M)=\max\{\,j:\ \wbar_j(TM)\neq0\,\}>0,
  \label{eq:delta}
\end{equation}
则 $\embdim(M)\ge d+\delta(M)$。
\end{theorem}

\begin{proof}[证明（个人复现）]
由 $TM\oplus\nu\cong\eps^{d+c}$ 得 $w(\nu)=w(TM)^{-1}$。$\nu$ 的秩为 $c$，故
$w_j(\nu)=0$（$j>c$）。于是 $\wbar_j(TM)=w_j(\nu)\neq0$ 迫使 $j\le c$。
\end{proof}

\begin{theorem}[欧氏嵌入的临界次数判据]\label{thm:critical}
设闭光滑 $d$ 维流形 $M$ 嵌入 $\R^{d+c}$ 且 $c>0$。则
\begin{equation}
  \wbar_i(TM)=0\qquad(i\ge c).
  \label{eq:critical}
\end{equation}
因而当 $\wbar_j(TM)\neq0$ 时，实际上有 $c\ge j+1$，环境维数至少为 $d+j+1$。
\end{theorem}

\begin{proof}
设 $\nu$ 为 $M$ 的秩 $c$ 法丛。$i>c$ 时 $w_i(\nu)=0$ 已由秩给出，故只需证
$w_c(\nu)=0$。

取向量 $u\in\R^{d+c}$，使对每个 $x\in M$ 都有 $u\notin T_xM$；这样的 $u$ 一定
存在，因为 $\bigcup_{x\in M}T_xM$ 是低维子空间的并，不能充满整个空间。把常向量场
$u$ 逐点正交投影到法向，得到 $\nu$ 的一个截面 $s(x)=u-(u\text{ 的切向分量})$。
由 $u\notin T_xM$ 知 $s(x)\neq0$ 处处成立，即 $\nu$ 有处处非零的截面。

对秩为 $c$ 的实向量丛，$w_c$ 就是它的模 $2$ Euler 类；而 Euler 类是零截面与任一
横截截面的相交数（模 $2$）。既然存在处处非零的截面，该相交数为零，故
$w_c(\nu)=0$。结合秩截断得 $w_i(\nu)=0$ 对一切 $i\ge c$ 成立，再由
$w(\nu)=\wbar(TM)$ 即得 \eqref{eq:critical}。

端点上的加强正是在这一步超出秩截断：秩截断只能控制 $i>c$，而处处非零截面的存在
额外给出 $i=c$ 的消失。Oproiu 在文献 \cite[p.~177]{Oproiu1977} 中引用这一判据，
并在第 181--184 页用它排除临界环境维数。
\end{proof}

\begin{remark}[两个定理的逻辑地位]
定理 \ref{thm:embedding-obstruction} 与定理 \ref{thm:critical} 都是嵌入的
\emph{必要}条件，而非充分条件：共轭类消失并不能构造嵌入。另一方面，非浸入的
下界可以转移为嵌入的下界（嵌入必为浸入），但浸入的上界、存在性与精确值都不作此
转移。
\end{remark}
